
\section{Interpretation and limitations}\label{sec:discussion}
The results support a specific methodological lesson: when evaluating reconstruction-derived features, membership alone is an inadequate control for the local-position information already available to a readout. Adding local coordinates produces the largest gain in this study. Reconstruction departures provide a further conditional benefit, but the matched RBF comparison does not support attributing that aggregate benefit to KAHM's inverse-Gram operation.

This conclusion is intentionally narrower than a ranking of general forecasting algorithms. The primary comparison uses a P-selected bank, one feature-scaling convention, one ridge family, four simulation laws, and three horizons. Changing the reference geometry, regularization, or readout could change the ordering. The linear control has a modest observed loss increase, and the nonlinear autoregression has an unfavorable K/P comparison. The tested representations are wider than their inputs. Neither the representation size nor the observed losses establish a generally efficient or robust forecasting method.

Additional experiments address related but different questions. A dimension-selected residual-forecast study establishes three of six primary criteria, but not its KAHM-over-RBF or forecasting-policy superiority criteria. A separate policy-augmentation experiment gives 0.816\% lower nonlinear loss than its base policy, without establishing the specified 3\% margin. On a descriptive 48-series M4 training-data panel, adding the KAHM candidate increases aggregate loss by 2.273\%. These findings are retained in the supplement, with their study-specific procedures and comparison families. They are not pooled with the 32 repetitions reported here. In particular, a useful feature block does not necessarily improve a policy that can already select other predictors.

The theoretical results also have a limited role. Equal raw spans permit different coefficient parameterizations, while distinct spans permit different functions. Neither statement orders regularized risks after the nonlinear transformations used in the experiment. A bounded radial map controls feature norm, and a standard centered-ridge argument bounds the magnitude of a fitted prediction conditional on its training sample. It does not supply an accuracy guarantee, a bound on reconstruction sensitivity near a cancelling denominator, or a stability guarantee for iterated forecasts.

Finally, the term membership describes local geometric weights, not a verified Markov state. We do not learn a transition operator for these weights or test closure of an observable space under dynamics. A direct forecast of $y_{t+h}$ cannot, by itself, establish autonomous representation dynamics. Such a claim would require a separate propagation model and evaluation of multistep predictions on a common target.

\section{Conclusion}
Within the specified simulation design, local position substantially improves forecasts from geometric memberships, and KAHM reconstruction departures further improve the position readout. A matched positive-kernel reconstruction achieves almost the same aggregate loss, so KAHM-specific superiority remains unestablished. Reference-span analysis explains why different reconstruction formulas may or may not add raw functions; the forecasting experiment shows why that algebra must be evaluated with a specified transformation, readout, and selection rule. The resulting contribution is a controlled account of when geometric enrichment is useful and of what the evidence does not establish.

\paragraph{Reproducibility and tool use.}
The supplementary material specifies all generators, numerical conventions, comparison families, and audit limitations. Machine-readable repetition losses support independent recomputation of the primary results. The accompanying repository documentation distinguishes document compilation, analysis of retained results, model replay, and simulation regeneration. AI tools assisted with protocol development, implementation, verification, analysis, and manuscript preparation. The authors are responsible for the scientific claims, evidence, and final text.
